\section{統計環境Rの導入}
\label{b03_introR}
\subsection{授業内容}
\subsubsection{概要}
携帯電話やタブレットが誰でも身近に利用できるようになってきてはいるが，それらはあくまでも受動的にサービスを利用する媒体であり，執筆や計算などユーザが能動的に電子機器を利用するにあたっては，パソコンの利用は不可欠である。
本講義だけでなく，心理学の研究をする上でも当然パソコンの利用は必須であり，とくに統計的データ処理をするにあたってはより深く知る必要がある。また，ExcelやNumbersなどの表計算でも記述統計量やグラフの作図はできるが，それ以上の「複数の変数を扱った多変量解析」や「群間の平均値の差を細かく検証する」といった心理統計的な作業には不十分なのである。
そこで，より専門的なツールとして統計環境Rを導入する。統計環境Rは専修大学人間科学部心理学科が公式に採用している統計環境であり，これを活用する総合開発環境であるRStudioの基本的な使い方，結果の再現性を担保するための工夫などについても言及する。それに先立って，コンピュータの基礎知識についても解説する。
\subsubsection{コマ主題細目}
\begin{description}

    \item[統計環境R/RStudioの準備] 心理学科では基本的に統計環境Rを共通のツールとして用いる。統計環境Rはフリーソフトウェアであり，誰でも無償で利用できる環境である。また，R単体では簡素な言語的やりとりができるに過ぎない。ファイルの操作や描画機能，記録，参照，管理など統計言語を扱うに当たって必要な総合的環境としてRStudioを併用することで，その利便性は飛躍的に向上する。まずはR/RStudioを導入し，基本的な画面構成やプロジェクトによる管理ができることを確認する。
	      \begin{flushright}
		      $\to$ Rについての入門書は多くあるがここでは\textcite{kosugi2019}のP.1--7を挙げておく。同じくRStudioを使った入門書は多くあるがここでは\textcite{kosugi2019}のP.7--24を挙げておく。
	      \end{flushright}

	\item[Rの基本操作]環境が整ったところで，Rの基本的な操作を見ていく。まずはRスクリプトをつかって，四則演算や関数計算，データの形式やパッケージの導入と利用について理解する。スクリプトペインとコンソールペインの違いを理解し，またスクリプトとして清書しておくことで計算手続きの記録がつけられることの利点を理解する。

	\item[基本的なデータ型と演算] RではNumeric(数値)，Character(文字列)，Logical(論理値)などの基本的なデータ型があり，それぞれの特徴を理解する。またベクトル，行列，データフレームといったデータ構造の違いを学び，基本的な統計量の算出方法を実践する。
          \begin{flushright}
              $\to$\textcite{kosugi2019}のP.25--40
          \end{flushright}

\end{description}
\subsubsection{キーワード}
\begin{itemize}
    \item R/RStudio
    \item スクリプトとコンソール
    \item データ型とデータ構造
    \item パッケージ管理
\end{itemize}
\subsection{授業情報}
\paragraph{コマの展開方法}
講義
\paragraph{標準シラバスにおける位置づけ}
科目番号5；心理学統計法；2.統計に関する基礎的な知識；C エクセル，R，SPSS入門(１) 統計分析のための言語の手ほどき　プログラムの基本的操作
\subsubsection{予習・復習課題}
\paragraph{予習}
事前にR(https://www.r-project.org/)とRStudio(https://posit.co/download/rstudio-desktop/)をダウンロードしてインストールしておくこと。インストールに問題がある場合は，授業開始前に教員やTAに相談すること。

\paragraph{復習}
授業で学んだR/RStudioの基本操作(四則演算，関数の使用，データの入力など)を実際に自分のPCで繰り返し練習すること。特に，スクリプトを用いた操作に慣れること。また，簡単な数値データを用いて平均値や標準偏差を算出する練習を行うこと。